\documentclass[10pt,a4paper]{amsart}
\usepackage[margin=19mm]{geometry}
\usepackage{amsmath,amssymb,mathtools}
\usepackage[scr=rsfs,cal=euler]{mathalfa}
\usepackage{xfrac}
\usepackage[unicode,hidelinks,bookmarks=false]{hyperref}
\newcommand{\ie}{i.e.\ }
\newcommand{\cf}{cf.\ }
\newcommand{\resp}{respectively\ }
\newcommand{\Subsection}{\subsection}
\providecommand{\nobreakdash}{\nobreak\hbox{-}}
\setlength{\emergencystretch}{2em}
\multlinegap0pt
\usepackage{bm}
\usepackage{bbm}
\usepackage{dsfont}
\usepackage{xspace}
\usepackage{cancel}
\usepackage{mathtools}
\usepackage{amsthm}
\makeatletter
\@ifundefined{Def}{\newtheorem{Def}{Def}}{}
\@ifundefined{Prop}{\newtheorem{Prop}[Def]{Proposition}}{}
\makeatother
\newcommand{\Area}{\operatorname{Area}}
\newcommand{\RR}{\mathbb{R}}
\newcommand{\Rot}{\operatorname{Rot}}
\newcommand{\SSp}{\mathbb{S}}
\newcommand{\cC}{\mathcal{C}}
\newcommand{\cL}{\mathcal{L}}
\newcommand{\eps}{\varepsilon}
\newcommand{\orig}{\mathbf{0}}
\newcommand{\sing}{\operatorname{sing}}
\title[Generic regularity for minimizing hypersurfaces in dimension]{Generic regularity for minimizing hypersurfaces in dimension 11}
\author{Otis Chodosh, Christos Mantoulidis, Felix Schulze, Zhihan Wang}
\date{Source: arXiv:2506.12852v1 (CC BY 4.0)}
\begin{document}
\maketitle

\noindent\emph{Source and licence.} Generic regularity for minimizing hypersurfaces in dimension 11. Version arXiv:2506.12852v1 (CC BY 4.0), distributed under the Creative Commons Attribution 4.0 International licence, as recorded in the arXiv licence metadata for this exact version and shown on the abstract page https://arxiv.org/abs/2506.12852v1. Full rights notice: licenses/arxiv-2506.12852-CC-BY-4.0.txt.\par\smallskip

\noindent\textit{Editorial scope.} Counted unit: the Prop whose source label is \texttt{prop:kappa.spine} in the article (source0.tex lines 1926--1929, source label \texttt{prop:kappa.spine}), delivered together with the article's own complete original proof of it (source0.tex lines 1967--2056, one \texttt{proof} environment of 90 lines). No mathematics is written, repaired or re-derived: statement and proof are the article's own text line for line. The proof is detached from the statement in the source: the article states the result at lines 1926--1929 and prints its proof at lines 1967--2056, and the proof environment's own header names the statement, so the association is the article's own. Required context retained uncounted: equ\_alpha(Omega) (lines 1947--1949); equ\_alpha\_exhaust limit (lines 1953--1955). Every definition, display and label of those lines is retained unchanged. Omitted from this package and not redistributed: 1--1925, 1930--1946, 1950--1952, 1956--1966, 2057--4208. Nothing inside a retained excerpt is omitted or altered: every retained body is delivered exactly as the article's lines are written, so no omission receipt of any kind accompanies this package. Citations: the retained excerpts carry no citation command at all, so no bibliography entry of the article is transcribed and no reference is redistributed. Reference adaptation: none was needed and none was made; every reference inside the delivered unit resolves inside it, so no unresolved reference or citation is produced. Printed slips kept literal: no printed word, symbol, hypothesis or number of the counted statement or of its proof was altered. Layout and class: the article's own class is \texttt{amsart} with packages including inputenc, geometry, enumerate, amsmath, amssymb, cite, amsfonts, mathrsfs, galois, appendix, xcolor, hyperref, marginnote, comment; the wrapper is the lane's pinned \texttt{amsart} editorial wrapper at 10pt on A4, which restates the theorem-like environments the retained text needs and adds the article's own macro definitions that the retained text needs (\texttt{\textbackslash Area}, \texttt{\textbackslash RR}, \texttt{\textbackslash Rot}, \texttt{\textbackslash SSp}, \texttt{\textbackslash cC}, \texttt{\textbackslash cL}, \texttt{\textbackslash eps}, \texttt{\textbackslash orig}, \texttt{\textbackslash sing}), with extra wrapper packages (bm, bbm, dsfont, xspace, cancel, mathtools). The delivered numbering is the unit's own. Assets: no retained span contains a figure environment, a graphics-inclusion command, a PDF-page inclusion, a table or a TikZ picture; no image file is copied into this package, the package declares no asset dependency and no image file is redistributed with it. Licence: version arXiv:2506.12852v1 of this article is distributed under the Creative Commons Attribution 4.0 International licence, as recorded in the arXiv licence metadata for this exact version and shown on the abstract page https://arxiv.org/abs/2506.12852v1. A full rights notice is delivered at licenses/arxiv-2506.12852-CC-BY-4.0.txt. Characters: every retained excerpt is pure ASCII, so the delivered engine, XeLaTeX with its default Latin Modern text font, reports no missing character.
  \begin{equation} \label{equ_alpha(Omega)}
	\alpha(\cC; \Omega) := \frac{n-2}{2} - \sqrt{\frac{(n-2)^2}{4} + \mu(\cC; \Omega)}.
  \end{equation}

  \begin{equation} \label{equ_alpha_exhaust limit}
    \lim_{j\to \infty}\alpha(\cC; \Omega_j) = \alpha(\cC; \cC\setminus Z) = \alpha(\cC)\,.       
  \end{equation}

\begin{Prop} \label{prop:kappa.spine}
	If $\cC \in \mathscr{C}_n \cap \Rot(\cC_\circ \times \RR^k)$, $\cC_\circ \in \mathscr{C}_{n-k}$, then
	\[ \alpha(\cC) = \alpha(\cC_\circ). \]
\end{Prop}

\begin{proof}[Proof of Proposition \ref{prop:kappa.spine}]
	Without loss of generality, we rotate $\RR^{n+1}$ so that
	\[ \cC = \cC_\circ \times \RR^k. \]
	Denote $n_\circ = n - k$, so that $\cC_\circ \subset \RR^{n_\circ + 1}$, and write $\cL$, $\cL_\circ$ for the links of $\cC$, $\cC_\circ$. 
	
    Fix an arbitrary domain $\Omega_\circ \Subset  \cL_\circ$ with smooth boundary. Denote
    \[ \Omega:= \SSp^n\cap (\operatorname{Cone}_\orig(\Omega_\circ)\times\RR^k)\subset \cL. \]
    The primary goal is to show that 
    \begin{equation} \label{equ_alpha(Omega)=alpha(Omega_0)}
      \alpha(\cC; \Omega) = \alpha(\cC_\circ; \Omega_\circ)\,. 
    \end{equation}
    It suffices to show this for $\Omega_\circ \neq \cL_\circ$.
    Note that combining \eqref{equ_alpha(Omega)=alpha(Omega_0)} and \eqref{equ_alpha_exhaust limit} proves Proposition \ref{prop:kappa.spine}. 

    To prove \eqref{equ_alpha(Omega)=alpha(Omega_0)} for $\Omega
    _\circ \neq \cL_\circ$, let $\mu_\circ:= \mu(\cC_\circ; \Omega_\circ)$ be the principal Dirichlet eigenvalue for $-\Delta_{\cL_\circ} - |A_{\cL_\circ}|^2$ on $\Omega_\circ$, $\varphi_\circ : \overline{\Omega_\circ} \to [0, \infty)$ be a corresponding principal Dirichlet eigenfunction, i.e.,
	\begin{equation} \label{eq:kappa.spine.phi}
		(\Delta_{\cL_\circ} + |A_{\cL_\circ}|^2) \varphi_\circ = -\mu_{\circ} \varphi_\circ \text{ in } \Omega_\circ, \quad \varphi_\circ = 0 \text{ on } \partial \Omega_\circ.
	\end{equation}
    Note that by strong maximum principle, $\varphi_\circ>0$ on $\Omega_\circ$. Also denote for simplicity $\alpha_\circ:= \alpha(\cC_\circ; \Omega_\circ)$. At this point, the assumption $\Omega_\circ\neq \cL_\circ$ enters to give $\alpha_\circ< \alpha(\cC_\circ)\leq (n_\circ-2)/2$. Let \[
      \psi: \operatorname{Cone}_\orig(\Omega_\circ)\times \RR^k \to [0, +\infty)\,, \quad \psi := r^{-\alpha_\circ}\varphi_\circ
    \] 
    be the $\RR^k$-invariant $(-\alpha_\circ)$-homogeneous extension of $\varphi_\circ$. Then 
	\begin{align*}
		 (\Delta_\cC + |A_\cC|^2)\psi 
        & = (\Delta_{\cC_\circ} + |A_{\cC_\circ}|^2) (r^{-\alpha_\circ} \varphi_\circ) \\
		& = (r^{-2} \Delta_{\cL_\circ} + (n_\circ-1) r^{-1} \partial_r + \partial_r^2 + r^{-2} |A_{\cL_\circ}|^2)(r^{-\alpha_\circ} \varphi_\circ) \\
		& = r^{-2} (\Delta_{\cL_\circ} + |A_{\cL_\circ}|^2) (r^{-\alpha_\circ} \varphi_\circ) + r^{-2} (\alpha_\circ^2 - (n_\circ-2) \alpha_\circ) (r^{-\alpha_\circ} \varphi_\circ) \\
		& = r^{-2} (-\mu_\circ) (r^{-\alpha_\circ} \varphi_\circ) + r^{-2} (\alpha_\circ^2 - (n_\circ-2) \alpha_\circ) (r^{-\alpha_{\circ}} \varphi_\circ) \\
		& = 0,
	\end{align*}
	where we used \eqref{eq:kappa.spine.phi}, \eqref{equ_alpha(Omega)} in the last two steps. Thus, $\psi$ is a positive Jacobi field with Dirichlet boundary conditions on $\operatorname{Cone}_\orig(\Omega_\circ) \times \RR^k \subset \cC$. Note, also, that $\psi$ being $(-\alpha_{\circ})$-homogeneous implies,
	\begin{equation} \label{eq:kappa.spine.homogeneous}
		\partial_\rho \psi = - \alpha_{\circ}\rho^{-1} \psi,
	\end{equation}
	where $\rho$ denotes the radial distance function from the origin in $\RR^{n+1}$. 
	
	Following the same computation as before but backwards, and in $\RR^{n+1}$ rather than $\RR^{n_\circ+1}$, we deduce from the Jacobi field nature of $\psi$ that:
	\begin{align*}
		0   & = (\Delta_{\cC} + |A_{\cC}|^2) \psi \\
			& = \left(\rho^{-2} \Delta_{\cL} + (n-1) \rho^{-1} \partial_\rho + \partial_\rho^2 + \rho^{-2} |A_{\cL}|^2\right) \psi \\
			& = \rho^{-2} (\Delta_{\cL} + |A_{\cL}|^2 + (\alpha_{\circ}^2 - (n-2) \alpha_{\circ})) \psi,
	\end{align*}
	where we used \eqref{eq:kappa.spine.homogeneous} in the last step. This implies that $\varphi:= \psi|_{\Omega}$ is a positive solution to 
    \begin{equation} \label{equ_varphi J_L eigenfunc}
      -(\Delta_{\cL} + |A_{\cL}|^2)\varphi = (\alpha_{\circ}^2 - (n-2) \alpha_{\circ})) \varphi \quad \text{ on }\Omega\,.         
    \end{equation}
    Hence, for every $\xi\in C^\infty_c(\Omega)$, let $\zeta:= \varphi^{-1}\xi$, we have 
    \begin{align*}
      \int_\Omega |\nabla_\cL\xi|^2 - |A_\cL|^2\xi^2 
      & = \int_\Omega |\nabla_\cL(\varphi\zeta)|^2 - |A_\cL|^2(\varphi\zeta)^2 \\
      & = \int_\Omega |\nabla_\cL\zeta|^2\varphi^2 + |\nabla_\cL\varphi|^2\zeta^2 + \varphi\nabla_\cL\varphi\cdot \nabla_\cL(\zeta^2) - |A_\cL|^2(\varphi\zeta)^2 \\
      & = \int_\Omega |\nabla_\cL\zeta|^2\varphi^2 - \zeta^2\varphi\cdot (\Delta_\cL + |A_\cL|^2)\varphi \\
      & \geq (\alpha_\circ^2 - (n-2)\alpha_\circ)\|\xi\|_{L^2(\Omega)}^2 \,.
    \end{align*}
    where we use \eqref{equ_varphi J_L eigenfunc} in the last inequality. This means \[
      \mu(\cC; \Omega) \geq \alpha_\circ^2 - (n-2)\alpha_\circ\,.
    \]
    Since by \eqref{equ_alpha(Omega)}, $\mu(\cC; \Omega) = \alpha(\cC; \Omega)^2 - (n-2)\alpha(\cC; \Omega)$ and $\alpha(\cC;\Omega)\leq (n-2)/2$, $\alpha_\circ< (n_\circ-2)/2<(n-2)/2$, this proves ``$\leq$" of \eqref{equ_alpha(Omega)=alpha(Omega_0)}.

    To prove the reverse inequality, we again use \eqref{equ_varphi J_L eigenfunc} to choose $\varphi$ as a test function after some cut-off near $\sing \cL = \bar \cL \setminus \cL$. Recall that by the definition of $\Omega$, we have \[
      \overline{\Omega}\cap \sing \cL = \SSp^n \cap (\RR^k\times \{\orig\}) \,.
    \]
    For every $\eps\in (0, 1)$, let $\eta_\eps:= \eta(r/\eps)$, viewed as a smooth function on $\cL$, where $\eta\in C^\infty(\RR; [0, 1])$ such that $\eta = 0$ on $(-\infty, 1/2]$, $\eta = 1$ on $[1, +\infty)$ and $|\eta'|\leq 10$.  
    Then $\eta_\eps\varphi$ is a Lipschitz function on $\cL$ supported away from $\sing \cL$ and vanishes outside $\Omega$, therefore by the same calculation as above, 
    \begin{align*}
      \mu(\cC; \Omega) \int_{\Omega}  (\eta_\eps\varphi)^2 
      & \leq \int_\Omega |\nabla_\cL(\eta_\eps\varphi)|^2 - |A_\cL|^2(\eta_\eps\varphi)^2 \\
      & = \int_\Omega |\nabla_\cL\eta_\eps|^2\varphi^2 - \eta_\eps^2\varphi\cdot (\Delta_\cL + |A_\cL|^2)\varphi \\
      & = \int_\Omega |\nabla_\cL\eta_\eps|^2\varphi^2 + (\alpha_\circ- (n-2)\alpha_\circ)\eta_\eps^2\varphi^2 \,.
    \end{align*}
    Recall that $\varphi = r^{-\alpha_\circ}\varphi_\circ$, where $\alpha_\circ<(n_\circ-2)/2$, we then have 
    \begin{align*}
      \int_{\Omega}(1-\eta_\eps^2)\varphi^2  
      & \leq \sum_{j=1}^\infty\int_{\Omega\cap \{2^{-j}\eps\leq r\leq 2^{-j+1}\eps\}} \|\varphi_\circ\|_{L^\infty}^2\cdot r^{-2\alpha_\circ} \\
      & \leq C(n, \varphi_\circ)\sum_{j=1}^\infty(2^{-j}\eps)^{-2\alpha_\circ} \Area(\cL \cap \{r<2^{-j+1}\eps\}).         
    \end{align*}
    Also,
    \begin{align*}
      \int_{\Omega} |\nabla_\cL\eta_\eps|^2\varphi^2 
        & \leq \int_{\Omega\cap \{\eps/2\leq r\leq\eps\}} \eps^{-2}|\eta'|^2\|\varphi_\circ\|_{L^\infty}^2\cdot r^{-2\alpha_\circ} \\
        & \leq C(n, \varphi_\circ)\eps^{-2-2\alpha_\circ} \Area(\cL \cap \{r<\eps\}).        
    \end{align*}
    By the monotonicity formula and a Vitali covering argument, \[
      \|\cL\|(\{r<\eps\})\leq C(\cL)\eps^{n_\circ} \,.  \]
    Combining these estimates and sending $\eps\to 0$, we derive, \[
      \mu(\cC; \Omega) \leq \alpha_\circ^2 - (n-2)\alpha_\circ \,.
    \]
    This proves ``$\geq$" of \eqref{equ_alpha(Omega)=alpha(Omega_0)}.
\end{proof}

\end{document}
